\section{Conclusiones}

El rediseño del portal institucional de la Universidad representó una nueva experiencia muy enriquecedora, no solo por la cantidad de información que maneja la escuela, sino por la oportunidad de modernizar la forma en que los usuarios interactúan con el sitio. Al finalizar este proyecto, se comprobó que utilizar una metodología basada en prototipos fue la decisión más acertada, ya que nos permitió avanzar sobre seguro. Hacer las cosas por pasos y revisar los diseños antes de escribir el código final ayudó a evitar funciones innecesarias que luego tendrían que ser descartadas, demostrando lo importante que es planear bien la estructura de un sistema web.

Además, trabajar de la mano con el departamento de Diseño Gráfico sirvió para entender cómo se dividen las responsabilidades en un entorno de desarrollo real. Fue muy interesante ver el flujo de trabajo y cómo un mismo proyecto pasa por procesos muy distintos: por un lado, la creatividad gráfica y por el otro, la lógica del código. Esta colaboración dejó muy claro que siempre debe existir un estándar a seguir; por ejemplo, la obligación de respetar los colores institucionales y mantener una misma línea visual en todos los elementos. Así, mientras ellos se encargaban de actualizar y modernizar el aspecto de ciertas secciones clave, el área de desarrollo se aseguraba de que esos nuevos diseños funcionaran perfectamente en la pantalla.

En la parte técnica, pasar de las tecnologías antiguas a herramientas más actuales fue un cambio muy necesario para la escuela. Más allá del cambio estético, el nuevo código estructurado resulta mucho más limpio, ordenado y escalable. Esto significa que el equipo interno de la UAT podrá darle mantenimiento o agregar nuevas funcionalidades de manera mucho más rápida y sencilla en el futuro. Participar en esta actualización me sirvió bastante para entender cómo se aplican los temas vistos en clase dentro de un proyecto real y la importancia de escribir buen código desde el principio.

Finalmente, tal y como se mostró en el capítulo de Resultados con las capturas del sistema terminado, el portal ya programado es la prueba directa de que se cumplieron los requerimientos iniciales del proyecto. Cada sección que se armó, desde el buscador general y el directorio hasta el mapa interactivo del campus, tiene el único propósito de facilitar el acceso a la información de la universidad para todo el público. El resultado es una página mucho más limpia e intuitiva, donde cualquier persona puede encontrar rápidamente lo que busca sin perderse.

\subsection{Trabajo Futuro}

Aunque el rediseño del portal cumple con todos los objetivos trazados para esta etapa, el ciclo de vida del software siempre permite áreas de mejora. Como trabajo futuro a corto y mediano plazo, se proponen las siguientes iniciativas:

\begin{itemize}
    \item \textbf{Pruebas con usuarios reales:} Realizar sesiones de prueba (\textit{testing}) directamente con estudiantes de nuevo ingreso para observar cómo navegan por la página y confirmar si encuentran la información de forma rápida.
    \item \textbf{Optimización de la vista móvil:} Dado que el sitio ya es responsivo, un siguiente paso sería pulir a detalle el diseño de ciertas secciones específicas. El objetivo es que la experiencia de lectura y navegación en teléfonos celulares sea todavía más fluida y cómoda para el usuario.
    \item \textbf{Mejoras en la navegación interna:} Conectar de mejor manera las distintas páginas del portal para hacer el recorrido del usuario mucho más intuitivo. En particular, se sugiere darle mayor visibilidad y enfoque a las áreas de Oferta Académica y Campus, ya que son de las más buscadas por los aspirantes.
    \item \textbf{Implementación de métricas de uso:} Agregar herramientas de analítica web para medir qué secciones son las más visitadas y desde qué dispositivos se conectan los alumnos. Obtener esta información serviría para seguir tomando decisiones y mejorando el portal basándose en datos reales.
\end{itemize}